% TOP_PROBLEM: {"id": 2806, "title": "Kirby Problem 3.8", "category_id": 7, "category": {"id": 7, "name": "topology", "display_name": "Topology", "description": "Properties preserved under continuous deformations.", "slug": "topology", "order_index": 7, "created_at": "2026-07-31T15:26:25.671Z"}, "status": "open", "problem_number": "KP-3.8", "set_id": 11}
% FIRST_POSTED: 2026-10-01
% ENTRY_KIND: statement_only
% UnsolvedMath ID 2806; source catalogue code KP-3.8
% !TeX program = lualatex
% Definition and sources only; this file is not an LLM solution attempt.
\documentclass[11pt,a4paper]{article}
\usepackage[margin=25mm]{geometry}
\usepackage{fontspec}
\setmainfont{lmroman10-regular.otf}[BoldFont=lmroman10-bold.otf,
 ItalicFont=lmroman10-italic.otf,BoldItalicFont=lmroman10-bolditalic.otf]
\usepackage{amsmath,amssymb,mathtools}
\usepackage{unicode-math}
\setmathfont{latinmodern-math.otf}
\AtBeginDocument{\let\setminus\smallsetminus}
\usepackage{xurl,hyperref}
\hypersetup{colorlinks=true,urlcolor=blue}
\setcounter{secnumdepth}{0}
\setlength{\parindent}{0pt}
\setlength{\parskip}{5pt}
\setlength{\emergencystretch}{3em}
\begin{document}
\section*{Kirby Problem 3.8}\label{2806}
{\small UnsolvedMath ID 2806 · KP-3.8 · Topology · Kirby's Problems in Low-Dimensional Topology}
\subsection{Definitions and mathematical statement}
For a group $G$, its profinite completion is the topological group
\[
 \widehat G=\varprojlim_{N\triangleleft G,\ [G:N]<\infty}G/N,
\]
where each finite quotient is discrete and the inverse limit carries its
compact inverse-limit topology. An isomorphism of profinite completions here
is an isomorphism of topological groups. This invariant records all finite
quotients together with their compatibility maps.

Let $M_1$ and $M_2$ be connected complete finite-volume hyperbolic
three-manifolds, with sectional curvature $-1$. Suppose
\[
 \widehat{\pi_1(M_1)}\cong\widehat{\pi_1(M_2)}.
\]
Must there be an isometry $M_1\to M_2$? Closed and cusped manifolds are both
included. Orientation reversal is allowed in the conclusion: the question
concerns their underlying hyperbolic manifolds rather than marked oriented ones.

No identification of the fundamental groups is part of the hypothesis.
The issue is whether finite quotient information alone forces that geometric
identification. Likewise, the question compares manifolds within the hyperbolic
class; it does not assert that every finitely generated group with the same
completion must itself be one of these manifold groups.
\subsection{Short English statement}
Does the complete collection of finite quotients of a hyperbolic three-manifold group determine the manifold up to isometry?
\subsection{Sources}
\begin{itemize}
\item[S1] ULAM.AI and the UnsolvedMath Contributors, supplied problems.json, record 2806 (KP-3.8), Kirby's Problems in Low-Dimensional Topology. Catalogue identity and source transcription; CC BY 4.0.
\url{https://huggingface.co/datasets/ulamai/UnsolvedMath}.
\item[S2] K3: A New Problem List in Low-Dimensional Topology, Problem 3.8. Expanded formulation of the supplied catalogue statement.
\url{https://math.berkeley.edu/sites/default/files/surv-295-ruberman-watermarked-author-pdf.pdf}.
\end{itemize}
\end{document}
